% OpenStax Calculus Volume 3 -- Section 2.7 Exercises: Cylindrical and Spherical Coordinates
% Exercises reproduced from OpenStax, licensed under CC BY 4.0.
% Source: https://openstax.org/books/calculus-volume-3/pages/2-7-cylindrical-and-spherical-coordinates
\documentclass{exercises}
\title{OpenStax Calculus Volume 3}
\subtitle{Section 2.7 Exercises: Cylindrical and Spherical Coordinates}
\sourceurl{https://openstax.org/books/calculus-volume-3/pages/2-7-cylindrical-and-spherical-coordinates}
\author{}
\date{}

\begin{document}
\maketitle


Use the following figure as an aid in identifying the relationship between the rectangular, cylindrical, and spherical coordinate systems.


\textit{[Figure: This figure is the first octant of the 3-dimensional coordinate system. There is a line segment from the origin upwards. It is labeled ``rho.'' The angle between this line segment and the $z$-axis is labeled ``phi.'' There is also a broken line from the origin to the shadow of the point. This line segment is in the $xy$-plane and is labeled ``r.'' The angle between $r$ and the $x$-axis is labeled ``theta.'']}


For the following exercises, the cylindrical coordinates $(r,\theta,z)$ of a point are given. Find the rectangular coordinates $(x,y,z)$ of the point.

\begin{enumerate}[start=1]
\item $(4,\frac{\pi}{6},3)$
\item $(3,\frac{\pi}{3},5)$
\item $(4,\frac{7\pi}{6},3)$
\item $(2,\pi,-4)$
\end{enumerate}

For the following exercises, the rectangular coordinates $(x,y,z)$ of a point are given. Find the cylindrical coordinates $(r,\theta,z)$ of the point.

\begin{enumerate}[resume]
\item $(1,\sqrt{3},2)$
\item $(1,1,5)$
\item $(3,-3,7)$
\item $(-2\sqrt{2},2\sqrt{2},4)$
\end{enumerate}

For the following exercises, the equation of a surface in cylindrical coordinates is given.


Find an equation of the surface in rectangular coordinates. Identify and graph the surface.

\begin{enumerate}[resume]
\item \techrequired $r=4$
\item \techrequired $z=r^{2}\cos^{2}\theta$
\item \techrequired $r^{2}\cos(2\theta)+z^{2}+1=0$
\item \techrequired $r=3\sin \theta$
\item \techrequired $r=2\cos \theta$
\item \techrequired $r^{2}+z^{2}=5$
\item \techrequired $r=2\sec \theta$
\item \techrequired $r=3\csc \theta$
\end{enumerate}

For the following exercises, the equation of a surface in rectangular coordinates is given. Find an equation of the surface in cylindrical coordinates.

\begin{enumerate}[resume]
\item $z=3$
\item $x=6$
\item $x^{2}+y^{2}+z^{2}=9$
\item $y=2x^{2}$
\item $x^{2}+y^{2}-16x=0$
\item $x^{2}+y^{2}-3\sqrt{x^{2}+y^{2}}+2=0$
\end{enumerate}

For the following exercises, the spherical coordinates $(\rho,\theta,\phi)$ of a point are given. Find the rectangular coordinates $(x,y,z)$ of the point.

\begin{enumerate}[resume]
\item $(3,0,\pi)$
\item $(1,\frac{\pi}{6},\frac{\pi}{6})$
\item $(12,-\frac{\pi}{4},\frac{\pi}{4})$
\item $(3,\frac{\pi}{4},\frac{\pi}{6})$
\end{enumerate}

For the following exercises, the rectangular coordinates $(x,y,z)$ of a point are given. Find the spherical coordinates $(\rho,\theta,\phi)$ of the point. Express the measure of the angles in degrees rounded to the nearest integer.

\begin{enumerate}[resume]
\item $(4,0,0)$
\item $(-1,2,1)$
\item $(0,3,0)$
\item $(-2,2\sqrt{3},4)$
\end{enumerate}

For the following exercises, the equation of a surface in spherical coordinates is given. Find an equation of the surface in rectangular coordinates. Identify and graph the surface.

\begin{enumerate}[resume]
\item \techrequired $\rho=3$
\item \techrequired $\phi=\frac{\pi}{3}$
\item \techrequired $\rho=2\cos \phi$
\item \techrequired $\rho=4\csc \phi$
\item \techrequired $\phi=\frac{\pi}{2}$
\item \techrequired $\rho=6\csc \phi \sec \theta$
\end{enumerate}

For the following exercises, the equation of a surface in rectangular coordinates is given. Find an equation of the surface in spherical coordinates. Identify the surface.

\begin{enumerate}[resume]
\item $x^{2}+y^{2}-3z^{2}=0$, $z\ne0$
\item $x^{2}+y^{2}+z^{2}-4z=0$
\item $z=6$
\item $x^{2}+y^{2}=9$
\end{enumerate}

For the following exercises, the cylindrical coordinates of a point are given. Find its associated spherical coordinates, with the measure of the angle $\phi$ in radians rounded to four decimal places.

\begin{enumerate}[resume]
\item \techrequired $(1,\frac{\pi}{4},3)$
\item \techrequired $(5,\pi,12)$
\item $(3,\frac{\pi}{2},3)$
\item $(3,-\frac{\pi}{6},3)$
\end{enumerate}

For the following exercises, the spherical coordinates of a point are given. Find its associated cylindrical coordinates.

\begin{enumerate}[resume]
\item $(2,-\frac{\pi}{4},\frac{\pi}{2})$
\item $(4,\frac{\pi}{4},\frac{\pi}{6})$
\item $(8,\frac{\pi}{3},\frac{\pi}{2})$
\item $(9,-\frac{\pi}{6},\frac{\pi}{3})$
\end{enumerate}

For the following exercises, find the most suitable system of coordinates to describe the solids.

\begin{enumerate}[resume]
\item The solid situated in the first octant with a vertex at the origin and enclosed by a cube of edge length $a$, where $a>0$
\item A spherical shell determined by the region between two concentric spheres centered at the origin, of radii of $a$ and $b$, respectively, where $b>a>0$
\item A solid inside sphere $x^{2}+y^{2}+z^{2}=9$ and outside cylinder ${(x-\frac{3}{2})}^{2}+y^{2}=\frac{9}{4}$
\item A cylindrical shell of height $10$ determined by the region between two cylinders with the same center, parallel rulings, and radii of $2$ and $5$, respectively
\item \techrequired Use a CAS to graph the region between elliptic paraboloid $z=x^{2}+y^{2}$ and cone $x^{2}+y^{2}-z^{2}=0$. Then describe the region in cylindrical coordinates.
\item \techrequired Use a CAS to graph in spherical coordinates the ``ice cream-cone region'' situated above the $xy$-plane between sphere $x^{2}+y^{2}+z^{2}=4$ and elliptical cone $x^{2}+y^{2}-z^{2}=0$.
\item Washington, DC, is located at $39^\circ \text{N}$ and $77^\circ \text{W}$ (see the following figure). Assume the radius of Earth is $4000$ mi. Express the location of Washington, DC, in spherical coordinates. \\ \textit{[Figure: This figure is an image of a globe. On the globe there is a point labeled where Washington, DC is located. It is labeled with 39 degrees north latitude and 77 degrees west longitude.]}
\item San Francisco is located at $37.78^\circ \text{N}$ and $122.42^\circ \text{W}$. Assume the radius of Earth is $4000$ mi. Express the location of San Francisco in spherical coordinates.
\item Find the latitude and longitude of Rio de Janeiro if its spherical coordinates are $(4000,-43.17^\circ,102.91^\circ)$.
\item Find the latitude and longitude of Berlin if its spherical coordinates are $(4000,13.38^\circ,37.48^\circ)$.
\item \techrequired Consider the torus of equation ${(x^{2}+y^{2}+z^{2}+R^{2}-r^{2})}^{2}=4R^{2}(x^{2}+y^{2})$, where $R\ge r>0$.
\begin{enumerate}[label=(\alph*)]
  \item Write the equation of the torus in spherical coordinates.
  \item If $R=r$, the surface is called a horn torus. Show that the equation of a horn torus in spherical coordinates is $\rho=2R\sin \phi$.
  \item Use a CAS to graph the horn torus with $R=r=2$ in spherical coordinates.
\end{enumerate}
\item \techrequired The ``bumpy sphere'' with an equation in spherical coordinates is $\rho=a+b\cos(m\theta)\sin(n\phi)$, with $\theta \in[0,2\pi]$ and $\phi \in[0,\pi]$, where $a$ and $b$ are positive numbers and $m$ and $n$ are positive integers, may be used in applied mathematics to model tumor growth.
\begin{enumerate}[label=(\alph*)]
  \item Show that the ``bumpy sphere'' is contained inside a sphere of equation $\rho=a+b$. Find the values of $\theta$ and $\phi$ at which the two surfaces intersect.
  \item Use a CAS to graph the surface for $a=14$, $b=2$, $m=4$, and $n=6$ along with sphere $\rho=a+b$.
  \item Find an equation of the intersection curve of the surface at b. with the cone $\phi=\frac{\pi}{12}$. Graph the intersection curve in the plane of intersection.
\end{enumerate}
\end{enumerate}

\end{document}
